\documentclass[11pt,a4paper]{article}
% main (standalone preamble for editing)
% vim: nowrap: spell:
% ══════════════════════════════════════════════════════════════════════════════
% Speed Maths --- shared preamble
% Included by every sheet and answer file via:  % main (standalone preamble for editing)
% vim: nowrap: spell:
% ══════════════════════════════════════════════════════════════════════════════
% Speed Maths --- shared preamble
% Included by every sheet and answer file via:  % main (standalone preamble for editing)
% vim: nowrap: spell:
% ══════════════════════════════════════════════════════════════════════════════
% Speed Maths --- shared preamble
% Included by every sheet and answer file via:  \input{../../shared/preamble}
% Editing THIS file changes the look of every sheet/answer across every pillar.
% ══════════════════════════════════════════════════════════════════════════════

\usepackage[a4paper, top=25mm, bottom=25mm, left=22mm, right=22mm]{geometry}
\usepackage{amsmath, amssymb}
\usepackage{enumitem}
\usepackage{booktabs}
\usepackage{array}
\usepackage{parskip}
\usepackage{microtype}
\usepackage{fancyhdr}
\usepackage{titlesec}
\usepackage{mdframed}
\usepackage{xcolor}
\usepackage[hidelinks]{hyperref}
\usepackage{graphicx}
\usepackage{tikz}
\usepackage{pgfplots}
\pgfplotsset{compat=1.18}

% ── Colours ───────────────────────────────────────────────────────────────────
\definecolor{darkgray}{gray}{0.15}
\definecolor{medgray}{gray}{0.45}
\definecolor{lightgray}{gray}{0.93}
\definecolor{boxgray}{gray}{0.88}

% ── Section formatting ──────────────────────────────────────────────────────────
\titleformat{\section}[block]
  {\normalfont\large\bfseries}{}
  {0pt}{}[\vspace{2pt}\hrule\vspace{6pt}]
\titlespacing*{\section}{0pt}{14pt}{4pt}

% ── List formatting ─────────────────────────────────────────────────────────────
\setlist[enumerate,1]{label=\textbf{\arabic*.}, leftmargin=2em, itemsep=10pt}

% ── Branding constants (edit here, not per-file) ────────────────────────────────
\newcommand{\SpeedExamLine}{\textbf{Speed Maths:} Competition \& Exam Prep}
\newcommand{\SpeedCredit}{Series by \href{https://www.linkedin.com/in/cerealdev/}{David Akinyele-Aje}}

% ── Header / footer ──────────────────────────────────────────────────────────────
% Usage (once, after \input{../../shared/preamble} and before \begin{document}):
%   \SpeedHeader{Algebra}{3}
% First arg  = pillar name shown as "Speed <Pillar> --- Daily Drill"
% Second arg = worksheet number
\pagestyle{fancy}
\newcommand{\SpeedHeader}[2]{%
  \fancyhf{}%
  \renewcommand{\headrulewidth}{0.4pt}%
  \setlength{\headheight}{14pt}%
  \fancyhead[L]{\small\textbf{Speed #1 --- Daily Drill}}%
  \fancyhead[R]{\small Worksheet \##2}%
  \fancyfoot[L]{\small \SpeedExamLine}%
  \fancyfoot[C]{\small\thepage}%
  \fancyfoot[R]{\small \SpeedCredit}%
  \renewcommand{\footrulewidth}{0.4pt}%
}

% ── Title block (used at the top of every sheet AND answer file) ───────────────
% Usage: \SpeedTitleBlock{Daily Algebra Drill \#3}{Jane Doe}
% %% AGENTS: When generating a new sheet, ask the human contributor for the name they want credited in argument 2.
% For answer files, pass the "--- Answers \& Investigations" suffix in the arg.
\newcommand{\SpeedTitleBlock}[2]{%
  \begin{center}
    {\LARGE\bfseries #1}\\[4pt]
    {\normalsize\itshape Sheet by #2}\\[6pt]
    \hrule\vspace{4pt}
    \begin{tabular}{llcc}
      \toprule
      \textbf{Section} & \textbf{Topic} & \textbf{Questions} & \textbf{Target time}\\
      \midrule
      A & Rapid Recognition          & 10 & 2 min 30 sec \\
      B & Manipulation Drills        & 10 & 8 min        \\
      C & Substitution \& Structure  & 8  & 10 min       \\
      D & Challenge                  & 5  & 15 min       \\
      \bottomrule
    \end{tabular}
    \vspace{4pt}
    \hrule
  \end{center}
}

% ── Sheet metadata: topic + the day's new tools ────────────────────────────────
% Usage (immediately after \SpeedTitleBlock, sheets only):
%   \SpeedMeta{Expanding, factorising, and the identity toolkit}
%             {difference of squares, sum/product form, completing the square}
%
% Arg 1 = the sheet's topic in one line. The website lists this instead of
%         "Sheet 03", so write the thing a reader would actually search for.
% Arg 2 = comma-separated, at most five: the tools this sheet introduces that
%         earlier sheets in the pillar did not. These become the website's tags.
%         Leave out anything the reader already met — the tags are meant to say
%         what is new today, not to summarise the sheet.
%
% Both are parsed straight out of this file by tools/build_website.py, so the
% sheet stays the single source of truth and the site cannot drift from it.
%
% This renders NOTHING. It is a declaration for the build script, not a piece of
% the sheet's layout: adding it to a sheet must not change that sheet's PDF, so
% the 25 existing sheets can be annotated without a single one of them being
% reissued. If the toolkit should also appear on the page, that is a separate
% decision about the sheet design — make it deliberately, here, not by leaving
% this macro printing something nobody asked it to.
\newcommand{\SpeedMeta}[2]{}

% ── Answer / Method / Investigate-further boxes (answer files only) ────────────
\newmdenv[
  backgroundcolor=lightgray,
  linecolor=medgray,
  linewidth=0.5pt,
  leftmargin=0pt, rightmargin=0pt,
  innerleftmargin=8pt, innerrightmargin=8pt,
  innertopmargin=5pt, innerbottommargin=5pt,
  skipabove=4pt, skipbelow=4pt
]{ansbox}

\newmdenv[
  backgroundcolor=white,
  linecolor=darkgray,
  linewidth=0.8pt,
  leftline=true, rightline=false, topline=false, bottomline=false,
  leftmargin=0pt, rightmargin=0pt,
  innerleftmargin=8pt, innerrightmargin=6pt,
  innertopmargin=4pt, innerbottommargin=4pt,
  skipabove=4pt, skipbelow=6pt
]{invbox}

\newcommand{\ans}[1]{%
  \begin{ansbox}
    \textbf{Answer:} #1
  \end{ansbox}}

\newcommand{\method}[1]{%
  {\small\textbf{Method:} #1}\par}

\newcommand{\inv}[1]{%
  \begin{invbox}
    \textbf{\small Investigate further:} {\small #1}
  \end{invbox}}

% ── Misc helpers ────────────────────────────────────────────────────────────────
\newcommand{\qn}[1]{\textbf{#1.}\;}

% Mistake-linking: attach a Top 5 Patterns / Common Traps bullet to the question
% label(s) in this sheet where it actually shows up, e.g. \seealso{B6, D2}
\newcommand{\seealso}[1]{\hfill\textit{\footnotesize(see #1)}}

% ── Graph options macro ─────────────────────────────────────────────────────────
% Usage: \GraphOptions{tikz code for A}{tikz code for B}{tikz code for C}{tikz code for D}
\newcommand{\GraphOptions}[4]{%
  \begin{center}
    \begin{tabular}{cc}
      \textbf{A)} & \textbf{B)} \\
      #1 & #2 \\[10pt]
      \textbf{C)} & \textbf{D)} \\
      #3 & #4
    \end{tabular}
  \end{center}
}

% ── Closing motivational rule + cross-promo footer (sheets only) ────────────────
% Usage: \SpeedClosing{"Quote text."}
\newcommand{\SpeedClosing}[1]{%
  \vspace{20pt}
  \begin{center}
  \rule{0.6\textwidth}{0.4pt}\\[4pt]
  {\small\itshape #1}\\[4pt]
  \rule{0.6\textwidth}{0.4pt}
  \end{center}
  \begin{center}
  {\small Found a mistake or want to contribute a sheet?}\\
  {\small Visit the open-source project at \href{https://github.com/The-CerealDev/speed-maths}{github.com/The-CerealDev/speed-maths}}
  \end{center}
}

% Editing THIS file changes the look of every sheet/answer across every pillar.
% ══════════════════════════════════════════════════════════════════════════════

\usepackage[a4paper, top=25mm, bottom=25mm, left=22mm, right=22mm]{geometry}
\usepackage{amsmath, amssymb}
\usepackage{enumitem}
\usepackage{booktabs}
\usepackage{array}
\usepackage{parskip}
\usepackage{microtype}
\usepackage{fancyhdr}
\usepackage{titlesec}
\usepackage{mdframed}
\usepackage{xcolor}
\usepackage[hidelinks]{hyperref}
\usepackage{graphicx}
\usepackage{tikz}
\usepackage{pgfplots}
\pgfplotsset{compat=1.18}

% ── Colours ───────────────────────────────────────────────────────────────────
\definecolor{darkgray}{gray}{0.15}
\definecolor{medgray}{gray}{0.45}
\definecolor{lightgray}{gray}{0.93}
\definecolor{boxgray}{gray}{0.88}

% ── Section formatting ──────────────────────────────────────────────────────────
\titleformat{\section}[block]
  {\normalfont\large\bfseries}{}
  {0pt}{}[\vspace{2pt}\hrule\vspace{6pt}]
\titlespacing*{\section}{0pt}{14pt}{4pt}

% ── List formatting ─────────────────────────────────────────────────────────────
\setlist[enumerate,1]{label=\textbf{\arabic*.}, leftmargin=2em, itemsep=10pt}

% ── Branding constants (edit here, not per-file) ────────────────────────────────
\newcommand{\SpeedExamLine}{\textbf{Speed Maths:} Competition \& Exam Prep}
\newcommand{\SpeedCredit}{Series by \href{https://www.linkedin.com/in/cerealdev/}{David Akinyele-Aje}}

% ── Header / footer ──────────────────────────────────────────────────────────────
% Usage (once, after % main (standalone preamble for editing)
% vim: nowrap: spell:
% ══════════════════════════════════════════════════════════════════════════════
% Speed Maths --- shared preamble
% Included by every sheet and answer file via:  \input{../../shared/preamble}
% Editing THIS file changes the look of every sheet/answer across every pillar.
% ══════════════════════════════════════════════════════════════════════════════

\usepackage[a4paper, top=25mm, bottom=25mm, left=22mm, right=22mm]{geometry}
\usepackage{amsmath, amssymb}
\usepackage{enumitem}
\usepackage{booktabs}
\usepackage{array}
\usepackage{parskip}
\usepackage{microtype}
\usepackage{fancyhdr}
\usepackage{titlesec}
\usepackage{mdframed}
\usepackage{xcolor}
\usepackage[hidelinks]{hyperref}
\usepackage{graphicx}
\usepackage{tikz}
\usepackage{pgfplots}
\pgfplotsset{compat=1.18}

% ── Colours ───────────────────────────────────────────────────────────────────
\definecolor{darkgray}{gray}{0.15}
\definecolor{medgray}{gray}{0.45}
\definecolor{lightgray}{gray}{0.93}
\definecolor{boxgray}{gray}{0.88}

% ── Section formatting ──────────────────────────────────────────────────────────
\titleformat{\section}[block]
  {\normalfont\large\bfseries}{}
  {0pt}{}[\vspace{2pt}\hrule\vspace{6pt}]
\titlespacing*{\section}{0pt}{14pt}{4pt}

% ── List formatting ─────────────────────────────────────────────────────────────
\setlist[enumerate,1]{label=\textbf{\arabic*.}, leftmargin=2em, itemsep=10pt}

% ── Branding constants (edit here, not per-file) ────────────────────────────────
\newcommand{\SpeedExamLine}{\textbf{Speed Maths:} Competition \& Exam Prep}
\newcommand{\SpeedCredit}{Series by \href{https://www.linkedin.com/in/cerealdev/}{David Akinyele-Aje}}

% ── Header / footer ──────────────────────────────────────────────────────────────
% Usage (once, after \input{../../shared/preamble} and before \begin{document}):
%   \SpeedHeader{Algebra}{3}
% First arg  = pillar name shown as "Speed <Pillar> --- Daily Drill"
% Second arg = worksheet number
\pagestyle{fancy}
\newcommand{\SpeedHeader}[2]{%
  \fancyhf{}%
  \renewcommand{\headrulewidth}{0.4pt}%
  \setlength{\headheight}{14pt}%
  \fancyhead[L]{\small\textbf{Speed #1 --- Daily Drill}}%
  \fancyhead[R]{\small Worksheet \##2}%
  \fancyfoot[L]{\small \SpeedExamLine}%
  \fancyfoot[C]{\small\thepage}%
  \fancyfoot[R]{\small \SpeedCredit}%
  \renewcommand{\footrulewidth}{0.4pt}%
}

% ── Title block (used at the top of every sheet AND answer file) ───────────────
% Usage: \SpeedTitleBlock{Daily Algebra Drill \#3}{Jane Doe}
% %% AGENTS: When generating a new sheet, ask the human contributor for the name they want credited in argument 2.
% For answer files, pass the "--- Answers \& Investigations" suffix in the arg.
\newcommand{\SpeedTitleBlock}[2]{%
  \begin{center}
    {\LARGE\bfseries #1}\\[4pt]
    {\normalsize\itshape Sheet by #2}\\[6pt]
    \hrule\vspace{4pt}
    \begin{tabular}{llcc}
      \toprule
      \textbf{Section} & \textbf{Topic} & \textbf{Questions} & \textbf{Target time}\\
      \midrule
      A & Rapid Recognition          & 10 & 2 min 30 sec \\
      B & Manipulation Drills        & 10 & 8 min        \\
      C & Substitution \& Structure  & 8  & 10 min       \\
      D & Challenge                  & 5  & 15 min       \\
      \bottomrule
    \end{tabular}
    \vspace{4pt}
    \hrule
  \end{center}
}

% ── Sheet metadata: topic + the day's new tools ────────────────────────────────
% Usage (immediately after \SpeedTitleBlock, sheets only):
%   \SpeedMeta{Expanding, factorising, and the identity toolkit}
%             {difference of squares, sum/product form, completing the square}
%
% Arg 1 = the sheet's topic in one line. The website lists this instead of
%         "Sheet 03", so write the thing a reader would actually search for.
% Arg 2 = comma-separated, at most five: the tools this sheet introduces that
%         earlier sheets in the pillar did not. These become the website's tags.
%         Leave out anything the reader already met — the tags are meant to say
%         what is new today, not to summarise the sheet.
%
% Both are parsed straight out of this file by tools/build_website.py, so the
% sheet stays the single source of truth and the site cannot drift from it.
%
% This renders NOTHING. It is a declaration for the build script, not a piece of
% the sheet's layout: adding it to a sheet must not change that sheet's PDF, so
% the 25 existing sheets can be annotated without a single one of them being
% reissued. If the toolkit should also appear on the page, that is a separate
% decision about the sheet design — make it deliberately, here, not by leaving
% this macro printing something nobody asked it to.
\newcommand{\SpeedMeta}[2]{}

% ── Answer / Method / Investigate-further boxes (answer files only) ────────────
\newmdenv[
  backgroundcolor=lightgray,
  linecolor=medgray,
  linewidth=0.5pt,
  leftmargin=0pt, rightmargin=0pt,
  innerleftmargin=8pt, innerrightmargin=8pt,
  innertopmargin=5pt, innerbottommargin=5pt,
  skipabove=4pt, skipbelow=4pt
]{ansbox}

\newmdenv[
  backgroundcolor=white,
  linecolor=darkgray,
  linewidth=0.8pt,
  leftline=true, rightline=false, topline=false, bottomline=false,
  leftmargin=0pt, rightmargin=0pt,
  innerleftmargin=8pt, innerrightmargin=6pt,
  innertopmargin=4pt, innerbottommargin=4pt,
  skipabove=4pt, skipbelow=6pt
]{invbox}

\newcommand{\ans}[1]{%
  \begin{ansbox}
    \textbf{Answer:} #1
  \end{ansbox}}

\newcommand{\method}[1]{%
  {\small\textbf{Method:} #1}\par}

\newcommand{\inv}[1]{%
  \begin{invbox}
    \textbf{\small Investigate further:} {\small #1}
  \end{invbox}}

% ── Misc helpers ────────────────────────────────────────────────────────────────
\newcommand{\qn}[1]{\textbf{#1.}\;}

% Mistake-linking: attach a Top 5 Patterns / Common Traps bullet to the question
% label(s) in this sheet where it actually shows up, e.g. \seealso{B6, D2}
\newcommand{\seealso}[1]{\hfill\textit{\footnotesize(see #1)}}

% ── Graph options macro ─────────────────────────────────────────────────────────
% Usage: \GraphOptions{tikz code for A}{tikz code for B}{tikz code for C}{tikz code for D}
\newcommand{\GraphOptions}[4]{%
  \begin{center}
    \begin{tabular}{cc}
      \textbf{A)} & \textbf{B)} \\
      #1 & #2 \\[10pt]
      \textbf{C)} & \textbf{D)} \\
      #3 & #4
    \end{tabular}
  \end{center}
}

% ── Closing motivational rule + cross-promo footer (sheets only) ────────────────
% Usage: \SpeedClosing{"Quote text."}
\newcommand{\SpeedClosing}[1]{%
  \vspace{20pt}
  \begin{center}
  \rule{0.6\textwidth}{0.4pt}\\[4pt]
  {\small\itshape #1}\\[4pt]
  \rule{0.6\textwidth}{0.4pt}
  \end{center}
  \begin{center}
  {\small Found a mistake or want to contribute a sheet?}\\
  {\small Visit the open-source project at \href{https://github.com/The-CerealDev/speed-maths}{github.com/The-CerealDev/speed-maths}}
  \end{center}
}
 and before \begin{document}):
%   \SpeedHeader{Algebra}{3}
% First arg  = pillar name shown as "Speed <Pillar> --- Daily Drill"
% Second arg = worksheet number
\pagestyle{fancy}
\newcommand{\SpeedHeader}[2]{%
  \fancyhf{}%
  \renewcommand{\headrulewidth}{0.4pt}%
  \setlength{\headheight}{14pt}%
  \fancyhead[L]{\small\textbf{Speed #1 --- Daily Drill}}%
  \fancyhead[R]{\small Worksheet \##2}%
  \fancyfoot[L]{\small \SpeedExamLine}%
  \fancyfoot[C]{\small\thepage}%
  \fancyfoot[R]{\small \SpeedCredit}%
  \renewcommand{\footrulewidth}{0.4pt}%
}

% ── Title block (used at the top of every sheet AND answer file) ───────────────
% Usage: \SpeedTitleBlock{Daily Algebra Drill \#3}{Jane Doe}
% %% AGENTS: When generating a new sheet, ask the human contributor for the name they want credited in argument 2.
% For answer files, pass the "--- Answers \& Investigations" suffix in the arg.
\newcommand{\SpeedTitleBlock}[2]{%
  \begin{center}
    {\LARGE\bfseries #1}\\[4pt]
    {\normalsize\itshape Sheet by #2}\\[6pt]
    \hrule\vspace{4pt}
    \begin{tabular}{llcc}
      \toprule
      \textbf{Section} & \textbf{Topic} & \textbf{Questions} & \textbf{Target time}\\
      \midrule
      A & Rapid Recognition          & 10 & 2 min 30 sec \\
      B & Manipulation Drills        & 10 & 8 min        \\
      C & Substitution \& Structure  & 8  & 10 min       \\
      D & Challenge                  & 5  & 15 min       \\
      \bottomrule
    \end{tabular}
    \vspace{4pt}
    \hrule
  \end{center}
}

% ── Sheet metadata: topic + the day's new tools ────────────────────────────────
% Usage (immediately after \SpeedTitleBlock, sheets only):
%   \SpeedMeta{Expanding, factorising, and the identity toolkit}
%             {difference of squares, sum/product form, completing the square}
%
% Arg 1 = the sheet's topic in one line. The website lists this instead of
%         "Sheet 03", so write the thing a reader would actually search for.
% Arg 2 = comma-separated, at most five: the tools this sheet introduces that
%         earlier sheets in the pillar did not. These become the website's tags.
%         Leave out anything the reader already met — the tags are meant to say
%         what is new today, not to summarise the sheet.
%
% Both are parsed straight out of this file by tools/build_website.py, so the
% sheet stays the single source of truth and the site cannot drift from it.
%
% This renders NOTHING. It is a declaration for the build script, not a piece of
% the sheet's layout: adding it to a sheet must not change that sheet's PDF, so
% the 25 existing sheets can be annotated without a single one of them being
% reissued. If the toolkit should also appear on the page, that is a separate
% decision about the sheet design — make it deliberately, here, not by leaving
% this macro printing something nobody asked it to.
\newcommand{\SpeedMeta}[2]{}

% ── Answer / Method / Investigate-further boxes (answer files only) ────────────
\newmdenv[
  backgroundcolor=lightgray,
  linecolor=medgray,
  linewidth=0.5pt,
  leftmargin=0pt, rightmargin=0pt,
  innerleftmargin=8pt, innerrightmargin=8pt,
  innertopmargin=5pt, innerbottommargin=5pt,
  skipabove=4pt, skipbelow=4pt
]{ansbox}

\newmdenv[
  backgroundcolor=white,
  linecolor=darkgray,
  linewidth=0.8pt,
  leftline=true, rightline=false, topline=false, bottomline=false,
  leftmargin=0pt, rightmargin=0pt,
  innerleftmargin=8pt, innerrightmargin=6pt,
  innertopmargin=4pt, innerbottommargin=4pt,
  skipabove=4pt, skipbelow=6pt
]{invbox}

\newcommand{\ans}[1]{%
  \begin{ansbox}
    \textbf{Answer:} #1
  \end{ansbox}}

\newcommand{\method}[1]{%
  {\small\textbf{Method:} #1}\par}

\newcommand{\inv}[1]{%
  \begin{invbox}
    \textbf{\small Investigate further:} {\small #1}
  \end{invbox}}

% ── Misc helpers ────────────────────────────────────────────────────────────────
\newcommand{\qn}[1]{\textbf{#1.}\;}

% Mistake-linking: attach a Top 5 Patterns / Common Traps bullet to the question
% label(s) in this sheet where it actually shows up, e.g. \seealso{B6, D2}
\newcommand{\seealso}[1]{\hfill\textit{\footnotesize(see #1)}}

% ── Graph options macro ─────────────────────────────────────────────────────────
% Usage: \GraphOptions{tikz code for A}{tikz code for B}{tikz code for C}{tikz code for D}
\newcommand{\GraphOptions}[4]{%
  \begin{center}
    \begin{tabular}{cc}
      \textbf{A)} & \textbf{B)} \\
      #1 & #2 \\[10pt]
      \textbf{C)} & \textbf{D)} \\
      #3 & #4
    \end{tabular}
  \end{center}
}

% ── Closing motivational rule + cross-promo footer (sheets only) ────────────────
% Usage: \SpeedClosing{"Quote text."}
\newcommand{\SpeedClosing}[1]{%
  \vspace{20pt}
  \begin{center}
  \rule{0.6\textwidth}{0.4pt}\\[4pt]
  {\small\itshape #1}\\[4pt]
  \rule{0.6\textwidth}{0.4pt}
  \end{center}
  \begin{center}
  {\small Found a mistake or want to contribute a sheet?}\\
  {\small Visit the open-source project at \href{https://github.com/The-CerealDev/speed-maths}{github.com/The-CerealDev/speed-maths}}
  \end{center}
}

% Editing THIS file changes the look of every sheet/answer across every pillar.
% ══════════════════════════════════════════════════════════════════════════════

\usepackage[a4paper, top=25mm, bottom=25mm, left=22mm, right=22mm]{geometry}
\usepackage{amsmath, amssymb}
\usepackage{enumitem}
\usepackage{booktabs}
\usepackage{array}
\usepackage{parskip}
\usepackage{microtype}
\usepackage{fancyhdr}
\usepackage{titlesec}
\usepackage{mdframed}
\usepackage{xcolor}
\usepackage[hidelinks]{hyperref}
\usepackage{graphicx}
\usepackage{tikz}
\usepackage{pgfplots}
\pgfplotsset{compat=1.18}

% ── Colours ───────────────────────────────────────────────────────────────────
\definecolor{darkgray}{gray}{0.15}
\definecolor{medgray}{gray}{0.45}
\definecolor{lightgray}{gray}{0.93}
\definecolor{boxgray}{gray}{0.88}

% ── Section formatting ──────────────────────────────────────────────────────────
\titleformat{\section}[block]
  {\normalfont\large\bfseries}{}
  {0pt}{}[\vspace{2pt}\hrule\vspace{6pt}]
\titlespacing*{\section}{0pt}{14pt}{4pt}

% ── List formatting ─────────────────────────────────────────────────────────────
\setlist[enumerate,1]{label=\textbf{\arabic*.}, leftmargin=2em, itemsep=10pt}

% ── Branding constants (edit here, not per-file) ────────────────────────────────
\newcommand{\SpeedExamLine}{\textbf{Speed Maths:} Competition \& Exam Prep}
\newcommand{\SpeedCredit}{Series by \href{https://www.linkedin.com/in/cerealdev/}{David Akinyele-Aje}}

% ── Header / footer ──────────────────────────────────────────────────────────────
% Usage (once, after % main (standalone preamble for editing)
% vim: nowrap: spell:
% ══════════════════════════════════════════════════════════════════════════════
% Speed Maths --- shared preamble
% Included by every sheet and answer file via:  % main (standalone preamble for editing)
% vim: nowrap: spell:
% ══════════════════════════════════════════════════════════════════════════════
% Speed Maths --- shared preamble
% Included by every sheet and answer file via:  \input{../../shared/preamble}
% Editing THIS file changes the look of every sheet/answer across every pillar.
% ══════════════════════════════════════════════════════════════════════════════

\usepackage[a4paper, top=25mm, bottom=25mm, left=22mm, right=22mm]{geometry}
\usepackage{amsmath, amssymb}
\usepackage{enumitem}
\usepackage{booktabs}
\usepackage{array}
\usepackage{parskip}
\usepackage{microtype}
\usepackage{fancyhdr}
\usepackage{titlesec}
\usepackage{mdframed}
\usepackage{xcolor}
\usepackage[hidelinks]{hyperref}
\usepackage{graphicx}
\usepackage{tikz}
\usepackage{pgfplots}
\pgfplotsset{compat=1.18}

% ── Colours ───────────────────────────────────────────────────────────────────
\definecolor{darkgray}{gray}{0.15}
\definecolor{medgray}{gray}{0.45}
\definecolor{lightgray}{gray}{0.93}
\definecolor{boxgray}{gray}{0.88}

% ── Section formatting ──────────────────────────────────────────────────────────
\titleformat{\section}[block]
  {\normalfont\large\bfseries}{}
  {0pt}{}[\vspace{2pt}\hrule\vspace{6pt}]
\titlespacing*{\section}{0pt}{14pt}{4pt}

% ── List formatting ─────────────────────────────────────────────────────────────
\setlist[enumerate,1]{label=\textbf{\arabic*.}, leftmargin=2em, itemsep=10pt}

% ── Branding constants (edit here, not per-file) ────────────────────────────────
\newcommand{\SpeedExamLine}{\textbf{Speed Maths:} Competition \& Exam Prep}
\newcommand{\SpeedCredit}{Series by \href{https://www.linkedin.com/in/cerealdev/}{David Akinyele-Aje}}

% ── Header / footer ──────────────────────────────────────────────────────────────
% Usage (once, after \input{../../shared/preamble} and before \begin{document}):
%   \SpeedHeader{Algebra}{3}
% First arg  = pillar name shown as "Speed <Pillar> --- Daily Drill"
% Second arg = worksheet number
\pagestyle{fancy}
\newcommand{\SpeedHeader}[2]{%
  \fancyhf{}%
  \renewcommand{\headrulewidth}{0.4pt}%
  \setlength{\headheight}{14pt}%
  \fancyhead[L]{\small\textbf{Speed #1 --- Daily Drill}}%
  \fancyhead[R]{\small Worksheet \##2}%
  \fancyfoot[L]{\small \SpeedExamLine}%
  \fancyfoot[C]{\small\thepage}%
  \fancyfoot[R]{\small \SpeedCredit}%
  \renewcommand{\footrulewidth}{0.4pt}%
}

% ── Title block (used at the top of every sheet AND answer file) ───────────────
% Usage: \SpeedTitleBlock{Daily Algebra Drill \#3}{Jane Doe}
% %% AGENTS: When generating a new sheet, ask the human contributor for the name they want credited in argument 2.
% For answer files, pass the "--- Answers \& Investigations" suffix in the arg.
\newcommand{\SpeedTitleBlock}[2]{%
  \begin{center}
    {\LARGE\bfseries #1}\\[4pt]
    {\normalsize\itshape Sheet by #2}\\[6pt]
    \hrule\vspace{4pt}
    \begin{tabular}{llcc}
      \toprule
      \textbf{Section} & \textbf{Topic} & \textbf{Questions} & \textbf{Target time}\\
      \midrule
      A & Rapid Recognition          & 10 & 2 min 30 sec \\
      B & Manipulation Drills        & 10 & 8 min        \\
      C & Substitution \& Structure  & 8  & 10 min       \\
      D & Challenge                  & 5  & 15 min       \\
      \bottomrule
    \end{tabular}
    \vspace{4pt}
    \hrule
  \end{center}
}

% ── Sheet metadata: topic + the day's new tools ────────────────────────────────
% Usage (immediately after \SpeedTitleBlock, sheets only):
%   \SpeedMeta{Expanding, factorising, and the identity toolkit}
%             {difference of squares, sum/product form, completing the square}
%
% Arg 1 = the sheet's topic in one line. The website lists this instead of
%         "Sheet 03", so write the thing a reader would actually search for.
% Arg 2 = comma-separated, at most five: the tools this sheet introduces that
%         earlier sheets in the pillar did not. These become the website's tags.
%         Leave out anything the reader already met — the tags are meant to say
%         what is new today, not to summarise the sheet.
%
% Both are parsed straight out of this file by tools/build_website.py, so the
% sheet stays the single source of truth and the site cannot drift from it.
%
% This renders NOTHING. It is a declaration for the build script, not a piece of
% the sheet's layout: adding it to a sheet must not change that sheet's PDF, so
% the 25 existing sheets can be annotated without a single one of them being
% reissued. If the toolkit should also appear on the page, that is a separate
% decision about the sheet design — make it deliberately, here, not by leaving
% this macro printing something nobody asked it to.
\newcommand{\SpeedMeta}[2]{}

% ── Answer / Method / Investigate-further boxes (answer files only) ────────────
\newmdenv[
  backgroundcolor=lightgray,
  linecolor=medgray,
  linewidth=0.5pt,
  leftmargin=0pt, rightmargin=0pt,
  innerleftmargin=8pt, innerrightmargin=8pt,
  innertopmargin=5pt, innerbottommargin=5pt,
  skipabove=4pt, skipbelow=4pt
]{ansbox}

\newmdenv[
  backgroundcolor=white,
  linecolor=darkgray,
  linewidth=0.8pt,
  leftline=true, rightline=false, topline=false, bottomline=false,
  leftmargin=0pt, rightmargin=0pt,
  innerleftmargin=8pt, innerrightmargin=6pt,
  innertopmargin=4pt, innerbottommargin=4pt,
  skipabove=4pt, skipbelow=6pt
]{invbox}

\newcommand{\ans}[1]{%
  \begin{ansbox}
    \textbf{Answer:} #1
  \end{ansbox}}

\newcommand{\method}[1]{%
  {\small\textbf{Method:} #1}\par}

\newcommand{\inv}[1]{%
  \begin{invbox}
    \textbf{\small Investigate further:} {\small #1}
  \end{invbox}}

% ── Misc helpers ────────────────────────────────────────────────────────────────
\newcommand{\qn}[1]{\textbf{#1.}\;}

% Mistake-linking: attach a Top 5 Patterns / Common Traps bullet to the question
% label(s) in this sheet where it actually shows up, e.g. \seealso{B6, D2}
\newcommand{\seealso}[1]{\hfill\textit{\footnotesize(see #1)}}

% ── Graph options macro ─────────────────────────────────────────────────────────
% Usage: \GraphOptions{tikz code for A}{tikz code for B}{tikz code for C}{tikz code for D}
\newcommand{\GraphOptions}[4]{%
  \begin{center}
    \begin{tabular}{cc}
      \textbf{A)} & \textbf{B)} \\
      #1 & #2 \\[10pt]
      \textbf{C)} & \textbf{D)} \\
      #3 & #4
    \end{tabular}
  \end{center}
}

% ── Closing motivational rule + cross-promo footer (sheets only) ────────────────
% Usage: \SpeedClosing{"Quote text."}
\newcommand{\SpeedClosing}[1]{%
  \vspace{20pt}
  \begin{center}
  \rule{0.6\textwidth}{0.4pt}\\[4pt]
  {\small\itshape #1}\\[4pt]
  \rule{0.6\textwidth}{0.4pt}
  \end{center}
  \begin{center}
  {\small Found a mistake or want to contribute a sheet?}\\
  {\small Visit the open-source project at \href{https://github.com/The-CerealDev/speed-maths}{github.com/The-CerealDev/speed-maths}}
  \end{center}
}

% Editing THIS file changes the look of every sheet/answer across every pillar.
% ══════════════════════════════════════════════════════════════════════════════

\usepackage[a4paper, top=25mm, bottom=25mm, left=22mm, right=22mm]{geometry}
\usepackage{amsmath, amssymb}
\usepackage{enumitem}
\usepackage{booktabs}
\usepackage{array}
\usepackage{parskip}
\usepackage{microtype}
\usepackage{fancyhdr}
\usepackage{titlesec}
\usepackage{mdframed}
\usepackage{xcolor}
\usepackage[hidelinks]{hyperref}
\usepackage{graphicx}
\usepackage{tikz}
\usepackage{pgfplots}
\pgfplotsset{compat=1.18}

% ── Colours ───────────────────────────────────────────────────────────────────
\definecolor{darkgray}{gray}{0.15}
\definecolor{medgray}{gray}{0.45}
\definecolor{lightgray}{gray}{0.93}
\definecolor{boxgray}{gray}{0.88}

% ── Section formatting ──────────────────────────────────────────────────────────
\titleformat{\section}[block]
  {\normalfont\large\bfseries}{}
  {0pt}{}[\vspace{2pt}\hrule\vspace{6pt}]
\titlespacing*{\section}{0pt}{14pt}{4pt}

% ── List formatting ─────────────────────────────────────────────────────────────
\setlist[enumerate,1]{label=\textbf{\arabic*.}, leftmargin=2em, itemsep=10pt}

% ── Branding constants (edit here, not per-file) ────────────────────────────────
\newcommand{\SpeedExamLine}{\textbf{Speed Maths:} Competition \& Exam Prep}
\newcommand{\SpeedCredit}{Series by \href{https://www.linkedin.com/in/cerealdev/}{David Akinyele-Aje}}

% ── Header / footer ──────────────────────────────────────────────────────────────
% Usage (once, after % main (standalone preamble for editing)
% vim: nowrap: spell:
% ══════════════════════════════════════════════════════════════════════════════
% Speed Maths --- shared preamble
% Included by every sheet and answer file via:  \input{../../shared/preamble}
% Editing THIS file changes the look of every sheet/answer across every pillar.
% ══════════════════════════════════════════════════════════════════════════════

\usepackage[a4paper, top=25mm, bottom=25mm, left=22mm, right=22mm]{geometry}
\usepackage{amsmath, amssymb}
\usepackage{enumitem}
\usepackage{booktabs}
\usepackage{array}
\usepackage{parskip}
\usepackage{microtype}
\usepackage{fancyhdr}
\usepackage{titlesec}
\usepackage{mdframed}
\usepackage{xcolor}
\usepackage[hidelinks]{hyperref}
\usepackage{graphicx}
\usepackage{tikz}
\usepackage{pgfplots}
\pgfplotsset{compat=1.18}

% ── Colours ───────────────────────────────────────────────────────────────────
\definecolor{darkgray}{gray}{0.15}
\definecolor{medgray}{gray}{0.45}
\definecolor{lightgray}{gray}{0.93}
\definecolor{boxgray}{gray}{0.88}

% ── Section formatting ──────────────────────────────────────────────────────────
\titleformat{\section}[block]
  {\normalfont\large\bfseries}{}
  {0pt}{}[\vspace{2pt}\hrule\vspace{6pt}]
\titlespacing*{\section}{0pt}{14pt}{4pt}

% ── List formatting ─────────────────────────────────────────────────────────────
\setlist[enumerate,1]{label=\textbf{\arabic*.}, leftmargin=2em, itemsep=10pt}

% ── Branding constants (edit here, not per-file) ────────────────────────────────
\newcommand{\SpeedExamLine}{\textbf{Speed Maths:} Competition \& Exam Prep}
\newcommand{\SpeedCredit}{Series by \href{https://www.linkedin.com/in/cerealdev/}{David Akinyele-Aje}}

% ── Header / footer ──────────────────────────────────────────────────────────────
% Usage (once, after \input{../../shared/preamble} and before \begin{document}):
%   \SpeedHeader{Algebra}{3}
% First arg  = pillar name shown as "Speed <Pillar> --- Daily Drill"
% Second arg = worksheet number
\pagestyle{fancy}
\newcommand{\SpeedHeader}[2]{%
  \fancyhf{}%
  \renewcommand{\headrulewidth}{0.4pt}%
  \setlength{\headheight}{14pt}%
  \fancyhead[L]{\small\textbf{Speed #1 --- Daily Drill}}%
  \fancyhead[R]{\small Worksheet \##2}%
  \fancyfoot[L]{\small \SpeedExamLine}%
  \fancyfoot[C]{\small\thepage}%
  \fancyfoot[R]{\small \SpeedCredit}%
  \renewcommand{\footrulewidth}{0.4pt}%
}

% ── Title block (used at the top of every sheet AND answer file) ───────────────
% Usage: \SpeedTitleBlock{Daily Algebra Drill \#3}{Jane Doe}
% %% AGENTS: When generating a new sheet, ask the human contributor for the name they want credited in argument 2.
% For answer files, pass the "--- Answers \& Investigations" suffix in the arg.
\newcommand{\SpeedTitleBlock}[2]{%
  \begin{center}
    {\LARGE\bfseries #1}\\[4pt]
    {\normalsize\itshape Sheet by #2}\\[6pt]
    \hrule\vspace{4pt}
    \begin{tabular}{llcc}
      \toprule
      \textbf{Section} & \textbf{Topic} & \textbf{Questions} & \textbf{Target time}\\
      \midrule
      A & Rapid Recognition          & 10 & 2 min 30 sec \\
      B & Manipulation Drills        & 10 & 8 min        \\
      C & Substitution \& Structure  & 8  & 10 min       \\
      D & Challenge                  & 5  & 15 min       \\
      \bottomrule
    \end{tabular}
    \vspace{4pt}
    \hrule
  \end{center}
}

% ── Sheet metadata: topic + the day's new tools ────────────────────────────────
% Usage (immediately after \SpeedTitleBlock, sheets only):
%   \SpeedMeta{Expanding, factorising, and the identity toolkit}
%             {difference of squares, sum/product form, completing the square}
%
% Arg 1 = the sheet's topic in one line. The website lists this instead of
%         "Sheet 03", so write the thing a reader would actually search for.
% Arg 2 = comma-separated, at most five: the tools this sheet introduces that
%         earlier sheets in the pillar did not. These become the website's tags.
%         Leave out anything the reader already met — the tags are meant to say
%         what is new today, not to summarise the sheet.
%
% Both are parsed straight out of this file by tools/build_website.py, so the
% sheet stays the single source of truth and the site cannot drift from it.
%
% This renders NOTHING. It is a declaration for the build script, not a piece of
% the sheet's layout: adding it to a sheet must not change that sheet's PDF, so
% the 25 existing sheets can be annotated without a single one of them being
% reissued. If the toolkit should also appear on the page, that is a separate
% decision about the sheet design — make it deliberately, here, not by leaving
% this macro printing something nobody asked it to.
\newcommand{\SpeedMeta}[2]{}

% ── Answer / Method / Investigate-further boxes (answer files only) ────────────
\newmdenv[
  backgroundcolor=lightgray,
  linecolor=medgray,
  linewidth=0.5pt,
  leftmargin=0pt, rightmargin=0pt,
  innerleftmargin=8pt, innerrightmargin=8pt,
  innertopmargin=5pt, innerbottommargin=5pt,
  skipabove=4pt, skipbelow=4pt
]{ansbox}

\newmdenv[
  backgroundcolor=white,
  linecolor=darkgray,
  linewidth=0.8pt,
  leftline=true, rightline=false, topline=false, bottomline=false,
  leftmargin=0pt, rightmargin=0pt,
  innerleftmargin=8pt, innerrightmargin=6pt,
  innertopmargin=4pt, innerbottommargin=4pt,
  skipabove=4pt, skipbelow=6pt
]{invbox}

\newcommand{\ans}[1]{%
  \begin{ansbox}
    \textbf{Answer:} #1
  \end{ansbox}}

\newcommand{\method}[1]{%
  {\small\textbf{Method:} #1}\par}

\newcommand{\inv}[1]{%
  \begin{invbox}
    \textbf{\small Investigate further:} {\small #1}
  \end{invbox}}

% ── Misc helpers ────────────────────────────────────────────────────────────────
\newcommand{\qn}[1]{\textbf{#1.}\;}

% Mistake-linking: attach a Top 5 Patterns / Common Traps bullet to the question
% label(s) in this sheet where it actually shows up, e.g. \seealso{B6, D2}
\newcommand{\seealso}[1]{\hfill\textit{\footnotesize(see #1)}}

% ── Graph options macro ─────────────────────────────────────────────────────────
% Usage: \GraphOptions{tikz code for A}{tikz code for B}{tikz code for C}{tikz code for D}
\newcommand{\GraphOptions}[4]{%
  \begin{center}
    \begin{tabular}{cc}
      \textbf{A)} & \textbf{B)} \\
      #1 & #2 \\[10pt]
      \textbf{C)} & \textbf{D)} \\
      #3 & #4
    \end{tabular}
  \end{center}
}

% ── Closing motivational rule + cross-promo footer (sheets only) ────────────────
% Usage: \SpeedClosing{"Quote text."}
\newcommand{\SpeedClosing}[1]{%
  \vspace{20pt}
  \begin{center}
  \rule{0.6\textwidth}{0.4pt}\\[4pt]
  {\small\itshape #1}\\[4pt]
  \rule{0.6\textwidth}{0.4pt}
  \end{center}
  \begin{center}
  {\small Found a mistake or want to contribute a sheet?}\\
  {\small Visit the open-source project at \href{https://github.com/The-CerealDev/speed-maths}{github.com/The-CerealDev/speed-maths}}
  \end{center}
}
 and before \begin{document}):
%   \SpeedHeader{Algebra}{3}
% First arg  = pillar name shown as "Speed <Pillar> --- Daily Drill"
% Second arg = worksheet number
\pagestyle{fancy}
\newcommand{\SpeedHeader}[2]{%
  \fancyhf{}%
  \renewcommand{\headrulewidth}{0.4pt}%
  \setlength{\headheight}{14pt}%
  \fancyhead[L]{\small\textbf{Speed #1 --- Daily Drill}}%
  \fancyhead[R]{\small Worksheet \##2}%
  \fancyfoot[L]{\small \SpeedExamLine}%
  \fancyfoot[C]{\small\thepage}%
  \fancyfoot[R]{\small \SpeedCredit}%
  \renewcommand{\footrulewidth}{0.4pt}%
}

% ── Title block (used at the top of every sheet AND answer file) ───────────────
% Usage: \SpeedTitleBlock{Daily Algebra Drill \#3}{Jane Doe}
% %% AGENTS: When generating a new sheet, ask the human contributor for the name they want credited in argument 2.
% For answer files, pass the "--- Answers \& Investigations" suffix in the arg.
\newcommand{\SpeedTitleBlock}[2]{%
  \begin{center}
    {\LARGE\bfseries #1}\\[4pt]
    {\normalsize\itshape Sheet by #2}\\[6pt]
    \hrule\vspace{4pt}
    \begin{tabular}{llcc}
      \toprule
      \textbf{Section} & \textbf{Topic} & \textbf{Questions} & \textbf{Target time}\\
      \midrule
      A & Rapid Recognition          & 10 & 2 min 30 sec \\
      B & Manipulation Drills        & 10 & 8 min        \\
      C & Substitution \& Structure  & 8  & 10 min       \\
      D & Challenge                  & 5  & 15 min       \\
      \bottomrule
    \end{tabular}
    \vspace{4pt}
    \hrule
  \end{center}
}

% ── Sheet metadata: topic + the day's new tools ────────────────────────────────
% Usage (immediately after \SpeedTitleBlock, sheets only):
%   \SpeedMeta{Expanding, factorising, and the identity toolkit}
%             {difference of squares, sum/product form, completing the square}
%
% Arg 1 = the sheet's topic in one line. The website lists this instead of
%         "Sheet 03", so write the thing a reader would actually search for.
% Arg 2 = comma-separated, at most five: the tools this sheet introduces that
%         earlier sheets in the pillar did not. These become the website's tags.
%         Leave out anything the reader already met — the tags are meant to say
%         what is new today, not to summarise the sheet.
%
% Both are parsed straight out of this file by tools/build_website.py, so the
% sheet stays the single source of truth and the site cannot drift from it.
%
% This renders NOTHING. It is a declaration for the build script, not a piece of
% the sheet's layout: adding it to a sheet must not change that sheet's PDF, so
% the 25 existing sheets can be annotated without a single one of them being
% reissued. If the toolkit should also appear on the page, that is a separate
% decision about the sheet design — make it deliberately, here, not by leaving
% this macro printing something nobody asked it to.
\newcommand{\SpeedMeta}[2]{}

% ── Answer / Method / Investigate-further boxes (answer files only) ────────────
\newmdenv[
  backgroundcolor=lightgray,
  linecolor=medgray,
  linewidth=0.5pt,
  leftmargin=0pt, rightmargin=0pt,
  innerleftmargin=8pt, innerrightmargin=8pt,
  innertopmargin=5pt, innerbottommargin=5pt,
  skipabove=4pt, skipbelow=4pt
]{ansbox}

\newmdenv[
  backgroundcolor=white,
  linecolor=darkgray,
  linewidth=0.8pt,
  leftline=true, rightline=false, topline=false, bottomline=false,
  leftmargin=0pt, rightmargin=0pt,
  innerleftmargin=8pt, innerrightmargin=6pt,
  innertopmargin=4pt, innerbottommargin=4pt,
  skipabove=4pt, skipbelow=6pt
]{invbox}

\newcommand{\ans}[1]{%
  \begin{ansbox}
    \textbf{Answer:} #1
  \end{ansbox}}

\newcommand{\method}[1]{%
  {\small\textbf{Method:} #1}\par}

\newcommand{\inv}[1]{%
  \begin{invbox}
    \textbf{\small Investigate further:} {\small #1}
  \end{invbox}}

% ── Misc helpers ────────────────────────────────────────────────────────────────
\newcommand{\qn}[1]{\textbf{#1.}\;}

% Mistake-linking: attach a Top 5 Patterns / Common Traps bullet to the question
% label(s) in this sheet where it actually shows up, e.g. \seealso{B6, D2}
\newcommand{\seealso}[1]{\hfill\textit{\footnotesize(see #1)}}

% ── Graph options macro ─────────────────────────────────────────────────────────
% Usage: \GraphOptions{tikz code for A}{tikz code for B}{tikz code for C}{tikz code for D}
\newcommand{\GraphOptions}[4]{%
  \begin{center}
    \begin{tabular}{cc}
      \textbf{A)} & \textbf{B)} \\
      #1 & #2 \\[10pt]
      \textbf{C)} & \textbf{D)} \\
      #3 & #4
    \end{tabular}
  \end{center}
}

% ── Closing motivational rule + cross-promo footer (sheets only) ────────────────
% Usage: \SpeedClosing{"Quote text."}
\newcommand{\SpeedClosing}[1]{%
  \vspace{20pt}
  \begin{center}
  \rule{0.6\textwidth}{0.4pt}\\[4pt]
  {\small\itshape #1}\\[4pt]
  \rule{0.6\textwidth}{0.4pt}
  \end{center}
  \begin{center}
  {\small Found a mistake or want to contribute a sheet?}\\
  {\small Visit the open-source project at \href{https://github.com/The-CerealDev/speed-maths}{github.com/The-CerealDev/speed-maths}}
  \end{center}
}
 and before \begin{document}):
%   \SpeedHeader{Algebra}{3}
% First arg  = pillar name shown as "Speed <Pillar> --- Daily Drill"
% Second arg = worksheet number
\pagestyle{fancy}
\newcommand{\SpeedHeader}[2]{%
  \fancyhf{}%
  \renewcommand{\headrulewidth}{0.4pt}%
  \setlength{\headheight}{14pt}%
  \fancyhead[L]{\small\textbf{Speed #1 --- Daily Drill}}%
  \fancyhead[R]{\small Worksheet \##2}%
  \fancyfoot[L]{\small \SpeedExamLine}%
  \fancyfoot[C]{\small\thepage}%
  \fancyfoot[R]{\small \SpeedCredit}%
  \renewcommand{\footrulewidth}{0.4pt}%
}

% ── Title block (used at the top of every sheet AND answer file) ───────────────
% Usage: \SpeedTitleBlock{Daily Algebra Drill \#3}{Jane Doe}
% %% AGENTS: When generating a new sheet, ask the human contributor for the name they want credited in argument 2.
% For answer files, pass the "--- Answers \& Investigations" suffix in the arg.
\newcommand{\SpeedTitleBlock}[2]{%
  \begin{center}
    {\LARGE\bfseries #1}\\[4pt]
    {\normalsize\itshape Sheet by #2}\\[6pt]
    \hrule\vspace{4pt}
    \begin{tabular}{llcc}
      \toprule
      \textbf{Section} & \textbf{Topic} & \textbf{Questions} & \textbf{Target time}\\
      \midrule
      A & Rapid Recognition          & 10 & 2 min 30 sec \\
      B & Manipulation Drills        & 10 & 8 min        \\
      C & Substitution \& Structure  & 8  & 10 min       \\
      D & Challenge                  & 5  & 15 min       \\
      \bottomrule
    \end{tabular}
    \vspace{4pt}
    \hrule
  \end{center}
}

% ── Sheet metadata: topic + the day's new tools ────────────────────────────────
% Usage (immediately after \SpeedTitleBlock, sheets only):
%   \SpeedMeta{Expanding, factorising, and the identity toolkit}
%             {difference of squares, sum/product form, completing the square}
%
% Arg 1 = the sheet's topic in one line. The website lists this instead of
%         "Sheet 03", so write the thing a reader would actually search for.
% Arg 2 = comma-separated, at most five: the tools this sheet introduces that
%         earlier sheets in the pillar did not. These become the website's tags.
%         Leave out anything the reader already met — the tags are meant to say
%         what is new today, not to summarise the sheet.
%
% Both are parsed straight out of this file by tools/build_website.py, so the
% sheet stays the single source of truth and the site cannot drift from it.
%
% This renders NOTHING. It is a declaration for the build script, not a piece of
% the sheet's layout: adding it to a sheet must not change that sheet's PDF, so
% the 25 existing sheets can be annotated without a single one of them being
% reissued. If the toolkit should also appear on the page, that is a separate
% decision about the sheet design — make it deliberately, here, not by leaving
% this macro printing something nobody asked it to.
\newcommand{\SpeedMeta}[2]{}

% ── Answer / Method / Investigate-further boxes (answer files only) ────────────
\newmdenv[
  backgroundcolor=lightgray,
  linecolor=medgray,
  linewidth=0.5pt,
  leftmargin=0pt, rightmargin=0pt,
  innerleftmargin=8pt, innerrightmargin=8pt,
  innertopmargin=5pt, innerbottommargin=5pt,
  skipabove=4pt, skipbelow=4pt
]{ansbox}

\newmdenv[
  backgroundcolor=white,
  linecolor=darkgray,
  linewidth=0.8pt,
  leftline=true, rightline=false, topline=false, bottomline=false,
  leftmargin=0pt, rightmargin=0pt,
  innerleftmargin=8pt, innerrightmargin=6pt,
  innertopmargin=4pt, innerbottommargin=4pt,
  skipabove=4pt, skipbelow=6pt
]{invbox}

\newcommand{\ans}[1]{%
  \begin{ansbox}
    \textbf{Answer:} #1
  \end{ansbox}}

\newcommand{\method}[1]{%
  {\small\textbf{Method:} #1}\par}

\newcommand{\inv}[1]{%
  \begin{invbox}
    \textbf{\small Investigate further:} {\small #1}
  \end{invbox}}

% ── Misc helpers ────────────────────────────────────────────────────────────────
\newcommand{\qn}[1]{\textbf{#1.}\;}

% Mistake-linking: attach a Top 5 Patterns / Common Traps bullet to the question
% label(s) in this sheet where it actually shows up, e.g. \seealso{B6, D2}
\newcommand{\seealso}[1]{\hfill\textit{\footnotesize(see #1)}}

% ── Graph options macro ─────────────────────────────────────────────────────────
% Usage: \GraphOptions{tikz code for A}{tikz code for B}{tikz code for C}{tikz code for D}
\newcommand{\GraphOptions}[4]{%
  \begin{center}
    \begin{tabular}{cc}
      \textbf{A)} & \textbf{B)} \\
      #1 & #2 \\[10pt]
      \textbf{C)} & \textbf{D)} \\
      #3 & #4
    \end{tabular}
  \end{center}
}

% ── Closing motivational rule + cross-promo footer (sheets only) ────────────────
% Usage: \SpeedClosing{"Quote text."}
\newcommand{\SpeedClosing}[1]{%
  \vspace{20pt}
  \begin{center}
  \rule{0.6\textwidth}{0.4pt}\\[4pt]
  {\small\itshape #1}\\[4pt]
  \rule{0.6\textwidth}{0.4pt}
  \end{center}
  \begin{center}
  {\small Found a mistake or want to contribute a sheet?}\\
  {\small Visit the open-source project at \href{https://github.com/The-CerealDev/speed-maths}{github.com/The-CerealDev/speed-maths}}
  \end{center}
}

\SpeedHeader{Sequences}{5}

\begin{document}

\SpeedTitleBlock{Daily Sequences Drill \#5 --- Answers \& Investigations}{David Akinyele-Aje}
\noindent\textit{New toolkit today: Monotonicity \& Boundedness $\cdot$ Fixed Points $\cdot$ Non-Modular Periodicity $\cdot$ Floor/Ceiling Sequences.}

%═══════════════════════════════════════════════════════════════════════════════
\section*{Section A \quad Rapid Recognition}
%═══════════════════════════════════════════════════════════════════════════════
\begin{enumerate}

%── A1 ──────────────────────────────────────────────────────────────────────────
\item Is the sequence $a_n=5-n$ increasing, decreasing, or neither?

\ans{Decreasing}
\method{$a_{n+1}-a_n=(5-(n+1))-(5-n)=-1<0$ for every $n$.}
\inv{Every AP with negative common difference is decreasing --- this is just Day 1's definition, restated in ``increasing/decreasing'' language.}

%── A2 ──────────────────────────────────────────────────────────────────────────
\item Is the sequence $a_n=n^2$ bounded above? Bounded below?

\ans{Not bounded above; bounded below (by $1$).}
\method{$n^2\to\infty$ as $n\to\infty$, so no upper bound exists. But $n^2\geq1$ for every positive integer $n$, so $1$ is a lower bound.}
\inv{A sequence can be bounded on one side only --- always check both directions separately.}

%── A3 ──────────────────────────────────────────────────────────────────────────
\item Find the fixed point of $f(x)=2x-3$ (i.e.\ solve $f(x)=x$).

\ans{$x=3$}
\method{$x=2x-3\Rightarrow-x=-3\Rightarrow x=3$.}
\inv{Check: $f(3)=2(3)-3=3$ \checkmark --- once a sequence generated by $f$ lands exactly on $3$, every later term stays at $3$ forever.}

%── A4 ──────────────────────────────────────────────────────────────────────────
\item A sequence satisfies $a_{n+1}=\dfrac1{a_n}$. If $a_1=2$, find $a_2,a_3,a_4$.

\ans{$a_2=\tfrac12,a_3=2,a_4=\tfrac12$}
\method{$a_2=\tfrac1{a_1}=\tfrac12$. $a_3=\tfrac1{a_2}=2$. $a_4=\tfrac1{a_3}=\tfrac12$.}
\inv{This sequence has period $2$: odd-indexed terms are $2$, even-indexed terms are $\tfrac12$, forever --- an early taste of today's periodicity theme.}

%── A5 ──────────────────────────────────────────────────────────────────────────
\item Compute $\lfloor3.7\rfloor$ and $\lceil3.2\rceil$.

\ans{$\lfloor3.7\rfloor=3$, $\lceil3.2\rceil=4$}
\method{Floor rounds down to the nearest integer; ceiling rounds up.}
\inv{Both operations leave an already-integer value unchanged --- $\lfloor5\rfloor=\lceil5\rceil=5$.}

%── A6 ──────────────────────────────────────────────────────────────────────────
\item Recall from Day 4 that a recurrence with both characteristic roots inside the unit circle converges to $0$. Does $a_n=\left(\tfrac12\right)^n+\left(\tfrac13\right)^n$ converge as $n\to\infty$? To what value?

\ans{Yes, converges to $0$.}
\method{Both $\left(\tfrac12\right)^n\to0$ and $\left(\tfrac13\right)^n\to0$ as $n\to\infty$, since both bases have magnitude less than $1$. Their sum also $\to0$.}
\inv{This is literally Day 4's D5 result, applied to specific numbers: $a_n=A r_1^n+Br_2^n$ with $A=B=1$, $r_1=\tfrac12,r_2=\tfrac13$, both inside the unit circle.}

%── A7 ──────────────────────────────────────────────────────────────────────────
\item Is the sequence $a_n=(-1)^n$ bounded? Does it converge?

\ans{Bounded (yes); does not converge.}
\method{Every term is either $1$ or $-1$, so $-1\leq a_n\leq1$ always (bounded). But the sequence never settles near a single value --- it alternates forever.}
\inv{This is the standard counterexample to ``boundedness implies convergence'' --- see A8 and Common Trap 1.}

%── A8 ──────────────────────────────────────────────────────────────────────────
\item True or false: every bounded sequence converges.

\ans{False}
\method{$a_n=(-1)^n$ (A7) is a direct counterexample: bounded, but not convergent.}
\inv{The correct combined statement is: bounded \emph{and monotonic} together imply convergence --- boundedness alone is not enough.}

%── A9 ──────────────────────────────────────────────────────────────────────────
\item Is the sequence $a_n=\dfrac n{n+1}$ increasing or decreasing? (Check $a_1,a_2,a_3$.)

\ans{Increasing}
\method{$a_1=\tfrac12,a_2=\tfrac23,a_3=\tfrac34$ --- each term is bigger than the last.}
\inv{B1 proves this rigorously for all $n$, not just the first three terms.}

%── A10 ─────────────────────────────────────────────────────────────────────────
\item Find the smallest $n$ for which $\lfloor\sqrt n\rfloor=4$.

\ans{$16$}
\method{$\lfloor\sqrt n\rfloor=4$ needs $4\leq\sqrt n<5$, i.e.\ $16\leq n<25$. The smallest such $n$ is $16$.}
\inv{This is the boundary case of C5's counting question --- always convert a floor equation into a double inequality on $n$ before searching.}

\end{enumerate}

%═══════════════════════════════════════════════════════════════════════════════
\section*{Section B \quad Manipulation Drills}
%═══════════════════════════════════════════════════════════════════════════════
\begin{enumerate}

%── B1 ──────────────────────────────────────────────────────────────────────────
\item Prove that $a_n=\dfrac n{n+1}$ is increasing for all $n\geq1$ by showing $a_{n+1}-a_n>0$.

\ans{Proof: $a_{n+1}-a_n=\dfrac1{(n+1)(n+2)}>0$.}
\method{$a_{n+1}-a_n=\dfrac{n+1}{n+2}-\dfrac n{n+1}=\dfrac{(n+1)^2-n(n+2)}{(n+2)(n+1)}=\dfrac{n^2+2n+1-n^2-2n}{(n+2)(n+1)}=\dfrac1{(n+1)(n+2)}$, which is positive for every $n\geq1$. $\square$}
\inv{This ``common denominator, then simplify the numerator'' technique for proving $a_{n+1}-a_n>0$ works on almost any rational sequence --- practice it on B2's boundedness question next.}

%── B2 ──────────────────────────────────────────────────────────────────────────
\item Is $a_n=\dfrac n{n+1}$ bounded above? If so, by what value (the supremum)?
\begin{itemize}
\item[A)] Yes, supremum $1$.
\item[B)] Yes, supremum $2$.
\item[C)] No, unbounded.
\item[D)] Yes, supremum $\tfrac12$.
\end{itemize}

\ans{A) Yes, supremum $1$.}
\method{$a_n=\dfrac n{n+1}=1-\dfrac1{n+1}<1$ for every $n$, and $a_n\to1$ as $n\to\infty$ --- so $1$ is the smallest possible upper bound (the supremum), even though no term ever equals it exactly.}
\inv{Rewriting $\dfrac n{n+1}$ as $1-\dfrac1{n+1}$ makes both the increasing behaviour (B1) and the bound (here) obvious at a glance --- always look for this kind of rewrite on a rational sequence.}

%── B3 ──────────────────────────────────────────────────────────────────────────
\item A sequence satisfies $a_{n+1}=\dfrac1{a_n}$ with $a_1=3$. Find $a_{100}$.
\begin{itemize}
\item[A)] $3$
\item[B)] $\tfrac13$
\item[C)] $100$
\item[D)] $\tfrac1{100}$
\end{itemize}

\ans{B) $\tfrac13$}
\method{The sequence has period $2$: $a_1=3,a_2=\tfrac13,a_3=3,a_4=\tfrac13,\ldots$ --- odd indices give $3$, even indices give $\tfrac13$. Since $100$ is even, $a_{100}=\tfrac13$.}
\inv{Always reduce a large index modulo the period before doing anything else --- $100\bmod2=0$, matching the ``even'' case.}

%── B4 ──────────────────────────────────────────────────────────────────────────
\item Find the fixed point(s) of $f(x)=\dfrac1x$ (i.e.\ solve $f(x)=x$).

\ans{$x=\pm1$}
\method{$\dfrac1x=x\Rightarrow x^2=1\Rightarrow x=\pm1$.}
\inv{A function can have more than one fixed point --- $f(x)=\tfrac1x$ has two, unlike A3's $f(x)=2x-3$, which has exactly one (any non-identity linear function has at most one).}

%── B5 ──────────────────────────────────────────────────────────────────────────
\item A sequence satisfies $a_{n+1}=f(a_n)$ where $f$ has a fixed point at $x=5$. If $a_1=5$, what is $a_n$ for every $n$?
\begin{itemize}
\item[A)] $a_n=5$ for all $n$.
\item[B)] $a_n=5n$.
\item[C)] $a_n\to5$ but never equals it.
\item[D)] Cannot be determined.
\end{itemize}

\ans{A) $a_n=5$ for all $n$.}
\method{$a_1=5$ is exactly the fixed point, so $a_2=f(5)=5$, and by the same reasoning every later term stays at $5$ forever.}
\inv{This is the simplest possible case of C8's general theorem: if a sequence \emph{starts} at a fixed point, it never leaves.}

%── B6 ──────────────────────────────────────────────────────────────────────────
\item Which of the following sequences is monotonic (always increasing or always decreasing)?
\begin{itemize}
\item[A)] $a_n=(-1)^n$
\item[B)] $a_n=\sin(n)$
\item[C)] $a_n=\dfrac1n$
\item[D)] $a_n=n\bmod3$
\end{itemize}

\ans{C) $\dfrac1n$}
\method{$\dfrac1n$ strictly decreases as $n$ increases. A) and B) oscillate without settling into one direction; D) cycles through $1,2,0,1,2,0,\ldots$ forever.}
\inv{A sequence is monotonic only if \emph{every} consecutive pair moves the same direction --- one exception anywhere disqualifies it, which is why periodic sequences (like D) are essentially never monotonic unless constant.}

%── B7 ──────────────────────────────────────────────────────────────────────────
\item Evaluate $\lfloor-2.3\rfloor$.
\begin{itemize}
\item[A)] $-2$
\item[B)] $-3$
\item[C)] $2$
\item[D)] $3$
\end{itemize}

\ans{B) $-3$}
\method{Floor always rounds toward $-\infty$, not toward $0$: the greatest integer that is still $\leq-2.3$ is $-3$, not $-2$.}
\inv{This is the single most common floor-function slip --- for negative numbers, ``round down'' means \emph{more} negative, not closer to zero.}

%── B8 ──────────────────────────────────────────────────────────────────────────
\item A sequence is defined by $a_n=\left\lfloor\dfrac n2\right\rfloor$. Find $a_1,a_2,a_3,a_4,a_5$.
\begin{itemize}
\item[A)] $0,1,1,2,2$
\item[B)] $1,1,2,2,3$
\item[C)] $0,0,1,1,2$
\item[D)] $0,1,2,2,3$
\end{itemize}

\ans{A) $0,1,1,2,2$}
\method{$a_1=\lfloor0.5\rfloor=0$, $a_2=\lfloor1\rfloor=1$, $a_3=\lfloor1.5\rfloor=1$, $a_4=\lfloor2\rfloor=2$, $a_5=\lfloor2.5\rfloor=2$.}
\inv{This sequence repeats each integer value exactly twice (except possibly the first) --- a simple example of how floor division ``compresses'' a sequence.}

%── B9 ──────────────────────────────────────────────────────────────────────────
\item Determine whether $a_n=\dfrac{(-1)^n}n$ converges, and if so, to what limit.

\ans{Converges to $0$.}
\method{$\left|\dfrac{(-1)^n}n\right|=\dfrac1n\to0$ as $n\to\infty$, and since $-\dfrac1n\leq a_n\leq\dfrac1n$, the sequence is squeezed to $0$ regardless of the alternating sign.}
\inv{This is a genuinely different situation from A7's $(-1)^n$: here the \emph{magnitude} shrinks to $0$ as well as the sign alternating, so the oscillation itself dies out and convergence does happen.}

%── B10 ─────────────────────────────────────────────────────────────────────────
\item For an AP with common difference $d$, under what condition on $d$ is the sequence strictly increasing?
\begin{itemize}
\item[A)] $d>0$
\item[B)] $d<0$
\item[C)] $d\neq0$
\item[D)] $d\geq0$
\end{itemize}

\ans{A) $d>0$}
\method{$a_{n+1}-a_n=d$ for every $n$ in an AP --- strictly increasing means this difference must be positive, i.e.\ $d>0$.}
\inv{Compare to Day 1: this is the same fact as Day 1's toolkit, now phrased in this sheet's ``increasing/decreasing'' language rather than ``common difference.''}

\end{enumerate}

%═══════════════════════════════════════════════════════════════════════════════
\section*{Section C \quad Substitution \& Structure}
%═══════════════════════════════════════════════════════════════════════════════
\begin{enumerate}

%── C1 ──────────────────────────────────────────────────────────────────────────
\item A sequence satisfies $a_{n+1}=\dfrac{a_n-1}{a_n+1}$, with $a_1=2$. Find $a_5$.
\begin{itemize}
\item[A)] $2$
\item[B)] $\tfrac13$
\item[C)] $-\tfrac12$
\item[D)] $-3$
\end{itemize}

\ans{A) $2$}
\method{$a_1=2$. $a_2=\dfrac{2-1}{2+1}=\dfrac13$. $a_3=\dfrac{1/3-1}{1/3+1}=\dfrac{-2/3}{4/3}=-\dfrac12$. $a_4=\dfrac{-1/2-1}{-1/2+1}=\dfrac{-3/2}{1/2}=-3$. $a_5=\dfrac{-3-1}{-3+1}=\dfrac{-4}{-2}=2$.}
\inv{This M\"obius map has exact period $4$: $a_5=a_1$, so the cycle $2,\tfrac13,-\tfrac12,-3$ repeats forever --- a classical fact about the transformation $x\mapsto\frac{x-1}{x+1}$.}

%── C2 ──────────────────────────────────────────────────────────────────────────
\item Continuing C1's sequence, what is $a_{101}$?
\begin{itemize}
\item[A)] $2$
\item[B)] $\tfrac13$
\item[C)] $-\tfrac12$
\item[D)] $-3$
\end{itemize}

\ans{A) $2$}
\method{The sequence has period $4$ (from C1). $101=4(25)+1$, so $101\equiv1\pmod4$, meaning $a_{101}=a_1=2$.}
\inv{Exactly the same index-reduction technique as B3, just with period $4$ instead of $2$ --- always find the period first, then reduce the target index modulo it.}

%── C3 ──────────────────────────────────────────────────────────────────────────
\item A sequence satisfies $a_1=1$ and $a_{n+1}=\sqrt{a_n+2}$. Which is true about $(a_n)$?
\begin{itemize}
\item[A)] Increasing and bounded above by $2$; converges to $2$.
\item[B)] Decreasing and bounded below by $1$.
\item[C)] Not monotonic; oscillates.
\item[D)] Unbounded, diverges to infinity.
\end{itemize}

\ans{A) Increasing and bounded above by $2$; converges to $2$.}
\method{Computing terms: $a_1=1,a_2=\sqrt3\approx1.73,a_3\approx1.93,a_4\approx1.98,\ldots$ --- visibly increasing and approaching $2$. (D1 proves this rigorously.)}
\inv{The fixed point of $f(x)=\sqrt{x+2}$ solves $x=\sqrt{x+2}\Rightarrow x^2=x+2\Rightarrow x^2-x-2=0\Rightarrow(x-2)(x+1)=0\Rightarrow x=2$ (rejecting $x=-1$ since $f(x)\geq0$ always) --- exactly matching the numerically observed limit.}

%── C4 ──────────────────────────────────────────────────────────────────────────
\item The first seven terms of a sequence of positive integers are: $w_1=10,w_2=15,w_3=18,w_4=22,w_5=31,w_6=36,w_7=41$. Consider the statement $(*)$: ``If $n$ is prime, then $w_n$ is a multiple of $3$ or a multiple of $5$.'' What is the smallest value of $n$ that provides a counterexample to $(*)$? \textit{\small(after TMUA 2021 Paper 2 Q10)}
\begin{itemize}
\item[A)] $2$
\item[B)] $3$
\item[C)] $5$
\item[D)] $7$
\end{itemize}

\ans{C) $5$}
\method{Check the primes in order: $n=2$: $w_2=15=3\times5$, a multiple of both --- not a counterexample. $n=3$: $w_3=18=3\times6$, a multiple of $3$ --- not a counterexample. $n=5$: $w_5=31$, which is neither a multiple of $3$ ($31=3\times10+1$) nor of $5$ --- this IS a counterexample. Since $5$ is the smallest prime that fails, it's the answer.}
\inv{The real TMUA question (2021 Paper 2 Q10) uses the sequence $15,21,30,37,44,51,59$ and the claim ``multiple of $3$ or $5$,'' also landing on $n=5$ as the smallest counterexample --- always check primes in increasing order and stop at the first failure, rather than checking every term first.}

%── C5 ──────────────────────────────────────────────────────────────────────────
\item A sequence is defined by $a_n=\lfloor\sqrt n\rfloor$. For how many values of $n$ in the range $1\leq n\leq100$ does $a_n=7$?
\begin{itemize}
\item[A)] $13$
\item[B)] $14$
\item[C)] $15$
\item[D)] $16$
\end{itemize}

\ans{C) $15$}
\method{$a_n=7$ needs $7\leq\sqrt n<8$, i.e.\ $49\leq n<64$. The integers from $49$ to $63$ inclusive number $63-49+1=15$.}
\inv{This is A10's boundary-finding technique extended to counting a whole range --- always convert to a double inequality on $n$ first, then count the integers in that range.}

%── C6 ──────────────────────────────────────────────────────────────────────────
\item A sequence is defined by $a_n=n-\lfloor n/3\rfloor\cdot3$ (the remainder when $n$ is divided by $3$, written using floor notation). Is $(a_n)$ periodic? If so, what is its period?
\begin{itemize}
\item[A)] Periodic, period $3$.
\item[B)] Periodic, period $2$.
\item[C)] Not periodic.
\item[D)] Periodic, period $1$ (constant).
\end{itemize}

\ans{A) Periodic, period $3$.}
\method{$a_n=n-3\lfloor n/3\rfloor$ is exactly the remainder of $n$ divided by $3$: for $n=1,2,3,4,5,6,\ldots$ this gives $1,2,0,1,2,0,\ldots$ --- repeating every $3$ terms.}
\inv{Floor notation and the $\bmod$ operator describe exactly the same thing here --- $n\bmod3$ and $n-3\lfloor n/3\rfloor$ are identical formulas written two different ways.}

%── C7 ──────────────────────────────────────────────────────────────────────────
\item A sequence satisfies $a_{n+1}=a_n^2-a_n+1$ for $n\geq1$, with $0<a_1<1$. Is $(a_n)$ increasing, decreasing, or could it be either depending on $a_1$?
\begin{itemize}
\item[A)] Always increasing (or constant), for any starting value.
\item[B)] Always decreasing.
\item[C)] Depends on $a_1$.
\item[D)] Cannot be determined without more information.
\end{itemize}

\ans{A) Always increasing (or constant), for any starting value.}
\method{$a_{n+1}-a_n=(a_n^2-a_n+1)-a_n=a_n^2-2a_n+1=(a_n-1)^2\geq0$ always, regardless of $a_n$'s value --- so $a_{n+1}\geq a_n$ for every $n$, for \emph{any} real starting value, not just $0<a_1<1$.}
\inv{The restriction ``$0<a_1<1$'' in the question is a red herring --- the algebra shows the result holds unconditionally. Recognising when a stated constraint isn't actually load-bearing is itself a skill (see Common Trap 5).}

%── C8 ──────────────────────────────────────────────────────────────────────────
\item For a sequence defined by $a_{n+1}=f(a_n)$ with $f$ continuous, if $(a_n)$ converges to a limit $L$, what must be true about $L$?
\begin{itemize}
\item[A)] $f(L)=L$ --- $L$ must be a fixed point of $f$.
\item[B)] $L=a_1$ always.
\item[C)] $f(L)=0$.
\item[D)] No general statement can be made.
\end{itemize}

\ans{A) $f(L)=L$ --- $L$ must be a fixed point of $f$.}
\method{Taking the limit of both sides of $a_{n+1}=f(a_n)$ as $n\to\infty$: the left side $\to L$, and since $f$ is continuous, the right side $\to f(L)$. So $L=f(L)$.}
\inv{This is the theorem underlying every fixed-point computation in this sheet (A3, B4, B5, C3, D1, D4) --- it's why solving $f(x)=x$ is always the right first move when a recursively-defined sequence is claimed (or suspected) to converge.}

\end{enumerate}

%═══════════════════════════════════════════════════════════════════════════════
\section*{Section D \quad Challenge}
%═══════════════════════════════════════════════════════════════════════════════
\begin{enumerate}

%── D1 ──────────────────────────────────────────────────────────────────────────
\item A sequence satisfies $a_1=1$ and $a_{n+1}=\sqrt{a_n+2}$ (as in C3). Prove rigorously, by induction, that $(a_n)$ is bounded above by $2$ for all $n$; then prove $(a_n)$ is strictly increasing; and hence explain why the sequence must converge, identifying its limit.

\ans{Bounded above by $2$, strictly increasing, converges to $2$.}
\method{\emph{Bounded above by $2$ (induction):} Base case $a_1=1<2$. Inductive step: assume $a_k<2$; then $a_{k+1}=\sqrt{a_k+2}<\sqrt{2+2}=2$. By induction, $a_n<2$ for all $n$.\\
\emph{Strictly increasing:} Since $0<a_n<2$, we have $(a_n-2)<0$ and $(a_n+1)>0$, so $(a_n-2)(a_n+1)<0$, i.e.\ $a_n^2-a_n-2<0$, i.e.\ $a_n+2>a_n^2$. Since $a_n\geq0$, taking square roots preserves the inequality: $\sqrt{a_n+2}>a_n$, i.e.\ $a_{n+1}>a_n$.\\
\emph{Convergence:} An increasing sequence bounded above must converge (Monotone Convergence Theorem). Taking the limit $L$ of $a_{n+1}=\sqrt{a_n+2}$: $L=\sqrt{L+2}\Rightarrow L^2=L+2\Rightarrow(L-2)(L+1)=0\Rightarrow L=2$ or $L=-1$. Since every $a_n>0$, $L\geq0$, so $L=2$. $\square$}
\inv{This is the complete, rigorous version of what C3 only observed numerically --- notice that all three pieces (boundedness, monotonicity, and the limit equation) reuse ideas from earlier in the sheet: induction, the $(a_n-2)(a_n+1)$ factoring trick from C7's style of algebra, and C8's fixed-point theorem.}

%── D2 ──────────────────────────────────────────────────────────────────────────
\item A M\"obius map $f(x)=\dfrac1{1-x}$ generates a sequence via $a_{n+1}=f(a_n)$. Starting from $a_1=2$, find the period of the resulting sequence (the smallest $k>0$ such that $a_{n+k}=a_n$ for all $n$).
\begin{itemize}
\item[A)] $2$
\item[B)] $3$
\item[C)] $4$
\item[D)] $6$
\end{itemize}

\ans{B) $3$}
\method{$a_1=2$. $a_2=\dfrac1{1-2}=-1$. $a_3=\dfrac1{1-(-1)}=\dfrac12$. $a_4=\dfrac1{1-\frac12}=2=a_1$.}
\inv{This is a different M\"obius map from C1/C2's --- $f(x)=\dfrac1{1-x}$ has order $3$ under composition (a classical fact), while C1's $f(x)=\dfrac{x-1}{x+1}$ has order $4$. Both are genuine, exact cycles --- always iterate directly rather than assuming a map is aperiodic.}

%── D3 ──────────────────────────────────────────────────────────────────────────
\item A sequence is defined by $a_1=100$ and $a_{n+1}=\lfloor a_n/2\rfloor$. Find the smallest $n$ for which $a_n=0$.
\begin{itemize}
\item[A)] $6$
\item[B)] $7$
\item[C)] $8$
\item[D)] $9$
\end{itemize}

\ans{C) $8$}
\method{$a_1=100,a_2=50,a_3=25,a_4=12,a_5=6,a_6=3,a_7=1,a_8=0$.}
\inv{Repeated floor-halving reaches $0$ in roughly $\log_2(a_1)$ steps --- here $\log_2100\approx6.6$, and the actual count ($7$ steps from $100$ to $0$, i.e.\ $a_8=0$) matches that estimate closely, since each floor operation can only lose at most one extra step compared to plain halving.}

%── D4 ──────────────────────────────────────────────────────────────────────────
\item A sequence satisfies $a_1=3$ and $a_{n+1}=\dfrac{a_n+\frac4{a_n}}2$. Which best describes its behaviour?
\begin{itemize}
\item[A)] Decreasing (after the first term) and bounded below by $2$; converges rapidly to $2$.
\item[B)] Increasing, converges to $2$.
\item[C)] Oscillates between values above and below $2$.
\item[D)] Diverges to infinity.
\end{itemize}

\ans{A) Decreasing (after the first term) and bounded below by $2$; converges rapidly to $2$.}
\method{$a_1=3$, $a_2=\dfrac{3+4/3}2=\dfrac{13/3}2=\dfrac{13}6\approx2.167$, $a_3\approx2.006$, $a_4\approx2.00002$ --- rapidly approaching $2$ from above. The fixed point solves $x=\dfrac{x+4/x}2\Rightarrow2x=x+\dfrac4x\Rightarrow x=\dfrac4x\Rightarrow x^2=4\Rightarrow x=2$ (taking the positive root, since all terms stay positive).}
\inv{This is Newton's method (the Babylonian method) for computing $\sqrt4$ --- the same iteration $a_{n+1}=\frac12\left(a_n+\frac c{a_n}\right)$ computes $\sqrt c$ for any $c>0$, converging extremely fast (the number of correct digits roughly doubles every step).}

%── D5 ──────────────────────────────────────────────────────────────────────────
\item A sequence satisfies $a_{n+1}=a_n(1-a_n)$ with $a_1=\tfrac12$. Compute $a_2,a_3$, and determine whether the sequence is monotonic.
\begin{itemize}
\item[A)] Decreasing for all $n$, since $a_{n+1}-a_n=-a_n^2\leq0$ always.
\item[B)] Increasing for all $n$.
\item[C)] Not monotonic.
\item[D)] Constant.
\end{itemize}

\ans{A) Decreasing for all $n$.}
\method{$a_2=\tfrac12\left(1-\tfrac12\right)=\tfrac14$. $a_3=\tfrac14\left(1-\tfrac14\right)=\tfrac14\cdot\tfrac34=\tfrac3{16}$. In general, $a_{n+1}-a_n=a_n(1-a_n)-a_n=a_n(1-a_n-1)=-a_n^2\leq0$ always, so the sequence never increases.}
\inv{This is C7's exact technique in reverse-flavour: there, $(a_n-1)^2\geq0$ forced the sequence up; here, $-a_n^2\leq0$ forces it down. Both times, rewriting $a_{n+1}-a_n$ as $\pm(\text{something})^2$ settles monotonicity instantly, without needing to compute many terms or invoke boundedness at all.}

\end{enumerate}

%═══════════════════════════════════════════════════════════════════════════════
\section*{Top 5 Patterns --- Worksheet \#5}
\begin{enumerate}[leftmargin=2em,itemsep=4pt]
  \item \textbf{Monotonicity via the sign of $a_{n+1}-a_n$.} Often collapses to $\pm(\text{something})^2$, settling the direction instantly without computing terms. \seealso{A1, A9, B1, C7, D5}
  \item \textbf{Fixed points: solve $f(L)=L$.} If a sequence converges, its limit must be a fixed point of the generating function. \seealso{A3, B4, B5, C8, D4}
  \item \textbf{Boundedness $+$ monotonicity together imply convergence.} Neither alone is enough --- see Common Trap 1. \seealso{C3, D1}
  \item \textbf{M\"obius maps can have small, exact periods under repeated composition.} Check by direct iteration --- don't assume a map is aperiodic just because it looks complicated. \seealso{C1, C2, D2}
  \item \textbf{Floor/ceiling notation is just an explicit ``round down/up.''} Expressions like $\lfloor n/k\rfloor$ often secretly encode periodicity (a remainder) or a counting bound. \seealso{A5, A10, B7, B8, C5, C6, D3}
\end{enumerate}

\section*{Common Traps --- Worksheet \#5}
\begin{enumerate}[leftmargin=2em,itemsep=4pt]
  \item \textbf{Assuming boundedness alone implies convergence.} $(-1)^n$ is bounded but never converges --- monotonicity is also needed. \seealso{A7, A8}
  \item \textbf{Forgetting floor rounds DOWN even for negative numbers.} $\lfloor-2.3\rfloor=-3$, not $-2$. \seealso{B7}
  \item \textbf{Assuming a sequence must be eventually constant or simply diverging.} Some sequences are genuinely, exactly periodic with period $>1$ forever. \seealso{C1, C2, D2}
  \item \textbf{Trusting a few computed terms as proof of monotonicity.} Numerical evidence is not a proof --- always confirm the sign of $a_{n+1}-a_n$ algebraically for all $n$. \seealso{D1}
  \item \textbf{Assuming a stated restriction on initial conditions is load-bearing.} Some monotonicity results hold far more generally than a question implies --- check whether the algebra actually needs the stated constraint. \seealso{C7}
\end{enumerate}

\SpeedClosing{``Solving $f(L)=L$ tells you where the limit would be, not that there is one.
The fixed point is a suspect, not a confession.''}

\end{document}
